

\section{Fourteenth Sunday after Trinity: With God there is no partiality}


\begin{quote}

Luke 4:23–30. And He said to them, You will surely say this proverb to Me: Physician, heal yourself! The great things which we have heard were done in Capernaum, do here also in Your own country! But He said, Truly I say to you, no prophet is well received in his own country. But in truth I say to you: There were many widows in Israel in the days of Elijah, when heaven was shut up for three years and six months, at the time when there was a great famine in the whole land, and to none of them was Elijah sent, except to Zarephath in Sidon, to a widow. And there were many lepers in Israel in the time of the prophet Elisha, and none of them was cleansed, except Naaman the Syrian. And all who were in the synagogue were filled with wrath when they heard this. And they rose up and thrust Him out of the town and led Him to the brow of the hill on which their town was built, in order to throw Him down. But He passed through the midst of them and went away.

\end{quote}

\bigskip

“There is no partiality with God” (Rom. 2:11): this is the great truth which Paul declares and maintains when, in opposition to the Jews’ claim upon salvation by virtue of their many advantages, he proclaims that a man is saved by faith alone, without regard to whether he is Jew or Greek.

It is the same thing Mary sings when Elizabeth greeted her as the Mother of God:

“He has shown strength with His arm; He has scattered those who are proud in the thought of their heart; He has cast down the mighty from their thrones and exalted the lowly; the hungry He has filled with good things, and the rich He has sent away empty” (Luke 1:51–53).

It is with this divine attribute as with the Word of God itself: it is two-edged; it always works either life or judgment, faith or offense. This blessed trait of God’s jealousy, that He does not regard persons, is the most blessed comfort of a poor sinner, and the bitterest offense of a rich righteous man.

This same divine attribute in Jesus was what at one and the same time saved the Canaanite woman, the leprous Samaritan, and the sighing publican, and stirred up the Jews’ bitterness, so that at last they hanged Him on the Cross and brought judgment down upon themselves.

In today’s text His own townsmen led Him out to throw Him down from a mountain.

And why?

He had come into their synagogue, had, as was His custom, read from the Word of God and spoken; He had chosen for His text the sixty-first chapter of Isaiah, where the prophet foretells the Messiah who shall come and proclaim “the gospel to the poor, heal those who have a broken heart, proclaim release to the captives, that the blind shall receive their sight, and that the oppressed shall receive their freedom.”

He had expounded these words concerning Himself and had spoken so blessedly that they all gave Him praise and marveled.

One would think that they would then also take the blessed words to heart, repent, and believe that this their townsman was in truth the Messiah, who could deliver them from the bondage of sin and give balm to broken hearts.

But no. No one was ever converted by finding a sermon excellent, still less by marveling that an acquaintance of theirs, whom they may formerly have held in high esteem, could preach so well.

A learned man once preached in a congregation, and spoke simply and movingly of sin and grace. Yes, they said, it was quite a good sermon, but it was no professor’s sermon. They had expected something different, something highly learned and striking, something that showed he had taken special pains for them; one felt disappointed, and there was half a spirit of dissatisfaction in the congregation.

The people of Nazareth also thought, no doubt, that these were blessed words Jesus had spoken to them; but they expected something more, something extraordinary, something specially intended for them, His own countrymen, some sign or wonder or great thing; and when they had to wait in vain, they were disappointed, bitterly disappointed. They forgot His blessed words and began to wonder by what right He could claim to be the one in whom Isaiah’s prophecy had been fulfilled—He?

“Is not this Joseph’s son?” they said.

Then it was Jesus, who knew the innermost thoughts of their hearts, who sharply cut it open and laid it before them. He revealed to them their frivolity and carnal curiosity, in that, over their desire to see wonders, they forgot the blessed Word, which is the power of God unto salvation for everyone who believes. He rebuked them for their proud vanity, according to which they considered themselves especially entitled to demand wonders from Him, their countryman, who had done such great things in Capernaum. And by examples drawn from the history of Israel He showed that the words which the Lord had spoken through Moses (Deut. 10:17), that there is no partiality with Him, He had acted upon and would act upon until the end of days. “He has mercy on whom He will, and whom He will He hardens” (Rom. 9:18).

Did the Lord look to the widow in Zarephath because she was a poor heathen woman, whom chosen Israel thought it had a right to despise? No, the Lord does not regard persons. He passed by the many widows in Israel and sent Elijah to this Phoenician widow, when she and her son were just about to die of hunger. While there was drought and famine everywhere else, even in the king’s house, her jar of meal and her cruse of oil never became empty so long as Elijah dwelt there; and Elijah raised her son again when he died. And why was such an abundance of blessing poured out over this despised woman? Because she believed Elijah’s word and trustfully baked a cake for him from the last meal she had. The Lord did not disappoint her childlike trust; for there is no partiality with Him. O, how blessed it is to cast all cares upon Him in the firm and unshaken faith that He is our jealous Father, who has care for us. But what was blessing for the widow in Zarephath, and salvation for every poor sinner who trustfully believes the Word of God, that was an offense to the inhabitants of Nazareth.

And what do you say of Naaman? He was a Syrian, an enemy of Israel, who even had a captive Jewish girl in his bondage. Nor was there any lack of lepers in Israel at that time; but Naaman the Syrian alone was healed by Elisha. And why? Was he less sinful than other men? O no!

When Elisha would have nothing to do with his gifts, paid no regard to his rank and dignity, but sent word that he was to wash himself seven times in the river Jordan, he became angry and cried out: “Are not Abana and Pharpar, the rivers of Syria, better than all the waters of Israel?” He was exactly as you and I have been when God has knocked upon our hearts and set before us His simple way of salvation. We have often taken offense at its simplicity. Had it been something difficult that God required of us, we would have been more willing. But only to believe in Jesus’ blood!

So it also went with Naaman. But little by little, as he looked upon his misery, his leprous body, and the prophet’s simple word was held before him, even by his servants, his proud heart yielded; he went down into the despised river Jordan, dipped himself under once, twice, three times—in vain; but in faith in the prophet’s word he persevered, and when he came up out of the water the seventh time, he was fresh and clean as a child. The Lord had saved him. He did not look to whether he was Syrian or Israelite, whether he was circumcised or uncircumcised; for there is no partiality with God.

But what was Naaman’s salvation, what was your salvation and mine when we came crawling to the Cross, that God did not, because of our leprosy, because of our many sins, cast us back, that same thing became the offense and judgment of the Nazarenes. When Jesus did not regard their persons and paid no heed to their lofty claims because He had been brought up among them, but on the contrary showed from God’s own Word that it is precisely such proud and unbelieving Pharisees whom God passes by with His jealousy, they became so enraged that they even laid hands on Him to kill Him. His hour had not yet come, and He went away through the midst of them. But they had laid bare the disposition of their hearts, the same disposition which in the blinded people of Israel brought forth the cry, Crucify, and the judgment that scattered them over the whole earth unto this day. They did not know the time of their visitation and took offense because God did not regard their persons and their hypocritical righteousness and the temple and circumcision.

And now, friend, what is this for you: There is no partiality with God? Has it become judgment for you or salvation? Has it become an aroma of life unto life or an aroma of death unto death?

Do you still have something to boast of? If not anything great, then at least something small? The Nazarenes could at least boast only that Jesus had been brought up in their midst. What have you? You have, after all, committed no gross sins; you are no murderer or whoremonger. Perhaps you belong to a congregation, an orthodox congregation—as it is so finely expressed—where there are even many God-fearing people and a pious shepherd of souls. Perhaps you read the Word of God diligently and go regularly to church and to the altar; you believe everything that stands in the Bible, and often under the sermon you have been moved to tears and to a new resolve. All this is good and well; but, friend, I tell you, remember Nazareth and fear! For there is no partiality with God, and tax collectors and harlots go into the kingdom of God before you.

Or perhaps you have nothing to boast of except your sin and your shame; you are leprous as Naaman and perishing from hunger as the widow in Zarephath; you are the chief of sinners and scarcely dare approach the Cross of Jesus, but stand far off and sigh: “God be merciful to me, a poor sinner!”

O, what a blessed message for you, that God does not regard persons; He will come into your wretched heart and let salvation come to your house, as He did with Zacchaeus. He will take away your leprosy and give you the Bread of Life, so that you shall neither hunger nor thirst for all eternity. But do you believe this?

Then come again with Him to Nazareth, listen again to Jesus’ reading and exposition of the sixty-first chapter in Isaiah, keep the blessed words in your heart, and lift up your soul in thanks and praise because “there is no partiality with God.”

